\section{The original integral theta reading and analytic norm estimates}
\label{sec:faithful-theta-analytic-norms}
The following source translations retain the actual integral model,
level data, vacuum and analytic-fibre identifications in their statements.
They prove the displayed reading, first-minimum, Gaussian and
projected-norm arguments under those same data. They do not identify an
arbitrary family with that model, and they do not supply the separate
analytic spanning theorem.


\section{Integer-exponent readings and bounded component multipliers}
\label{sec:sr-source-integer-reading}
This section follows the actual negative-exponent clearing argument,
rather than replacing it by an extension principle.
We distinguish the Fourier propositions used by the source's
bounded-section algebra from the algebraic section-reading argument.

\subsection{The bounded-section product}
Fix the source Hermitian frame datum \(X\), its open frame domain
\(\mathcal D_X\), and integral row weights \(\lambda,\lambda'\).
Write
\[
 \|P\|_{\lambda,x}=|P(x)|\,n_\lambda(x),\qquad
 n_\lambda(x)>0,\qquad
 n_{\lambda+\lambda'}(x)=n_\lambda(x)n_{\lambda'}(x).
 \tag{IR1}\label{eq:sr-source-reading-nfac}
\]
The last equality follows directly from the product of positive
successive Gram-minor ratios in \(n_\lambda\), using
\(r^{a+b}=r^ar^b\) for \(r>0\).
A bounded section means a holomorphic section with the source
Borel and arithmetic transformation laws and bounded weighted norm.
The source space \(SS_\lambda\) consists of such sections whose
zero Fourier coefficient vanishes along every admissible rational
isotropic direction.

Here are the two exact Fourier inputs used at this point:
\begin{enumerate}
\item A nonzero Fourier coefficient of a moderate holomorphic section
along an admissible isotropic direction has frequency
\(\nu=(\nu_\tau)\) with \(\nu_\tau\ge0\) for every \(\tau\).
A bounded section is moderate with exponent \(0\).
\item The coefficient of a product is the convergent convolution of
the two coefficient families:
\[
 \operatorname{fj}_{\lambda+\lambda'}(PQ;\xi,\nu,x)
 =\sum_{\nu'}\operatorname{fj}_\lambda(P;\xi,\nu',x)
       \operatorname{fj}_{\lambda'}(Q;\xi,\nu-\nu',x).
 \tag{IR2}\label{eq:sr-source-reading-convolution}
\]
The sum is the source's \(\mathrm{HasSum}\) assertion, not a
formal multiplication of series.
\end{enumerate}
These are exactly \srcid{ModerateCoefficientPositivity} and
\srcid{CoefficientConvolution}.

\begin{srlemma}\label{lem:sr-source-bdd-cusp-product}
Under these two Fourier inputs, \(P\in SS_\lambda\) and a bounded
section \(Q\) of weight \(\lambda'\) imply \(PQ\in SS_{\lambda+\lambda'}\).
\end{srlemma}
\begin{proof}
Holomorphy is preserved by multiplication.
Both transformation characters multiply; for the Borel character
this is the identity
\(\prod_{\tau,i}b_{\tau,ii}^{\lambda_{\tau,i}+\lambda'_{\tau,i}}
 =b_\lambda(b)b_{\lambda'}(b)\).
Every diagonal entry is nonzero.
Equation~\eqref{eq:sr-source-reading-nfac} gives
\(\|PQ\|_{\lambda+\lambda',x}
 =\|P\|_{\lambda,x}\|Q\|_{\lambda',x}\), so the product is bounded.

Apply \eqref{eq:sr-source-reading-convolution} at \(\nu=0\).
For a nonzero convolution term, the first input gives
\(\nu'_\tau\ge0\) and \(-\nu'_\tau\ge0\) for every \(\tau\).
Thus \(\nu'=0\), and the \(P\)-factor is zero by cuspidality.
Every term is therefore zero.
Uniqueness of the stated sum proves that the product's zero
coefficient vanishes. This is
\srcid{BoundedSectionAlgebra.siegelSpaceAlgebra_of}, at its
\(SS\)-multiplication clause.
\end{proof}
The same argument with the triangle inequality proves that \(SS\)
is closed under zero, addition and scalar multiplication.
In taking sums of coefficients, integrability on the finite-volume
period domain follows from boundedness: real isotropic flow
preserves all the Gram minors, hence \(n_\lambda\), and the
Fourier exponential has modulus one.
Consequently the source's span-defined \(\mathrm{holSub}\) equals
the space of families whose every component belongs to \(SS\).
This is \srcid{mem_holSub_model}.

\subsection{The actual integer-exponent reading}
Use the source \(E,H,D_w,l,X_o,c\) and an actual good complex
uniformization \(U\).
Suppose \(p+q\ge1\), a source key section
\(\kappa\in\mathrm{KeySec}(X_o,c)\) has been supplied, and
\(G=(G_1,G_2)\) has even \(G_2\) and antitone integer exponent
rows:
\[
 i\le j<\operatorname{mult}_{D_w}(\tau)
 \ \Longrightarrow\ G_1(\tau,j)\le G_1(\tau,i).
 \tag{IR3}\label{eq:sr-source-integer-antitone}
\]
Let \(f\) be an actual algebraic section in
\(\mathrm{secZ}(X_o,c,\mathrm{exZ}(D_w,G))\).
Retain the uniformization's actual reading law for dominant
natural weights and even tensor degree.
Also retain the actual nonvanishing section supplier:
for every \(x_0\in\mathcal D_{U,k}\) there are \(m\ge0\) and
\[
 f_0\in\mathrm{sec}(X_o,c,0,2(m+1)),\qquad
             \operatorname{rd}_U(f_0)_k(x_0)\ne0.
 \tag{IR4}\label{eq:sr-source-reading-nonvanishing}
\]
The source supplies (IR4) by its level-\(3\) eta rule in
\srcid{UnfoldedIsSecZeroKeySection.exists_sec_rdU_ne_zero}.

\begin{srproposition}\label{prop:sr-source-integer-reading}
For every component \(k\), the actual reading
\(\operatorname{rd}_U(f)_k\) is a section of the exact model weight
\(\operatorname{modelWt}(U,G,k)\).
\end{srproposition}
\begin{proof}
Fix \(x_0\) and choose (IR4).
Write \(G_2=2g_2\), set \(N=|G_2|\), and choose \(B\ge0\) with
\(-B\le G_1(\tau,j)\) for every actual row index.
The finite sum of all \(|G_1(\tau,j)|\) is one possible \(B\).
Set exactly as in the source
\[
 \begin{aligned}
 t_g&=(B+1)m+(N+1)(m+1),&t&=g_2+t_g,\\
 a_g(\tau,i)&=
 \begin{cases}B+1&i<\operatorname{mult}_{D_w}(\tau),\\0&\text{otherwise},
 \end{cases}&
 a(\tau,i)&=
 \begin{cases}G_1(\tau,i)+B+1&i<\operatorname{mult}_{D_w}(\tau),\\
 0&\text{otherwise}.
 \end{cases}
 \end{aligned}
 \tag{IR5}\label{eq:sr-source-integer-clearing}
\]
The lower bound on \(G_1\) makes \(a\) natural-valued.
Moreover \(t_g\ge N+1\), so \(t\ge0\).
Both \(a_g\) and \(a\) are antitone, including the transition
to the zero entries outside the actual rows.

Let \(\rho=\kappa.\rho\), the specified algebraic unit of key
exponent, and put
\[
                  g=\rho^{B+1}f_0^{B+N+2}.
 \tag{IR6}\label{eq:sr-source-reading-clearing-section}
\]
The exact exponent identities are
\[
 \begin{split}
 \mathrm{exZ}(D_w,\mathrm{wtZ}(a_g,2t_g))
  &=(B+1)\,\mathrm{keyEx}(D_w)
      +(B+N+2)\mathrm{exZ}(D_w,\mathrm{wtZ}(0,2(m+1))),\\
 \mathrm{exZ}(D_w,\mathrm{wtZ}(a,2t))
  &=\mathrm{exZ}(D_w,G)
       +\mathrm{exZ}(D_w,\mathrm{wtZ}(a_g,2t_g)).
 \end{split}
 \tag{IR7}\label{eq:sr-source-reading-exponent-identities}
\]
For the tensor coordinate the first identity is
\((B+1)(-2)+2(B+N+2)(m+1)=2t_g\);
the second uses \(G_2=2g_2\).
For every determinant coordinate they read respectively
\(B+1=B+1\) and
\(G_1+B+1=G_1+(B+1)\).
Multiplication and powers of actual algebraic sections therefore give
\[
       g\in\mathrm{sec}(a_g,2t_g),\qquad
       fg\in\mathrm{sec}(a,2t).
\]
The uniformization's given dominant-weight reading law makes their
readings sections of the corresponding natural model weights.

The model weights add, because both original tensor weights are
even and their algebraic exponents agree in (IR7).
Thus, writing \(\mu=\operatorname{modelWt}(U,G,k)\) and
\(\lambda=\mathrm{laU}(U,a_g,2t_g,k)\), the readings of
\(fg\) and \(g\) have weights \(\mu+\lambda\) and \(\lambda\).

Their denominator reading at \(x_0\) is nonzero.
Indeed \(\operatorname{rd}_U(f_0)_k(x_0)\ne0\).
Multiplicativity then forces the reading of \(1\) to be \(1\)
there. Since \(\rho\rho^{-1}=1\), the reading of \(\rho\) is
nonzero there too. Equation (IR6) now gives the assertion.
Finally,
\[
    \operatorname{rd}_U(f)_k\,\operatorname{rd}_U(g)_k
                     =\operatorname{rd}_U(fg)_k
\]
holds on the whole frame domain.

For completeness, the source's local quotient argument gives
all three section clauses. Near \(x_0\) the denominator stays
nonzero, so the quotient is holomorphic.
For an upper triangular invertible \(b\),
\(b_\lambda(b)\ne0\).
Evaluate the displayed identity at \(x_0\) and \(x_0b\) and
cancel \(b_\lambda(b)\operatorname{rd}_U(g)_k(x_0)\);
this gives the Borel weight \(\mu\).
For an arithmetic transformation \(\gamma\), the multiplier
\(\chi_\lambda(\gamma)\) is nonzero, since it has finite order.
The identical cancellation gives the multiplier
\(\chi_\mu(\gamma)\).
Since \(x_0\) was arbitrary, all section clauses hold everywhere.
These are precisely
\srcid{IsSecReadingDecreasingIntegerExponents.isSec_of_local_quot}
and \srcid{isSec_modelWt_rdU}.
\end{proof}

\subsection{Preservation of the actual component cusp space}
Let \(Z\) be the source component, and let \(s\) have the same
hypotheses as \(f\) in the preceding proposition, with weight \(G_0\).
Assume also
\[
 s\in\mathrm{suppOn}(Z),\qquad
 \forall\sigma\text{ good}\ \exists C_\sigma\
 \forall z\in\mathrm{framesOn}(U_\sigma,Z):\
                    \mathrm{homN}(s;z)\le C_\sigma .
 \tag{IR8}\label{eq:sr-source-component-bdd}
\]
All frames and norms here belong to the same supplied
uniformization. No uniformity in a varying family is being added.

\begin{srproposition}\label{prop:sr-source-component-cusp-multiplier}
Under the two Fourier inputs above, for every \(G\) with even
tensor coordinate,
\[
      f\in\mathrm{cuspOn}(U,Z,G)
          \quad\Longrightarrow\quad
      fs\in\mathrm{cuspOn}(U,Z,G+G_0).
 \tag{IR9}\label{eq:sr-source-component-cusp-product}
\]
\end{srproposition}
\begin{proof}
Algebraic section exponents add, and support of \(s\) on \(Z\)
gives support of \(fs\) there.
Fix a good \(\sigma\).
The actual model reading of \(s\) is a section by the preceding
proposition.
On \(\mathrm{framesOn}(U_\sigma,Z)\), its weighted norm is at most
\(C_\sigma/\mathrm{lieFacZ}(D_w,x,G_0)\), whose denominator is
strictly positive.
Off those component frames its reading is zero by the support
condition. It is therefore a bounded section on every component,
with bound
\[
                \max\{0,C_\sigma/\mathrm{lieFacZ}(D_w,x,G_0)\}.
\]
The component readings of \(f\) belong to \(SS\).
The first lemma puts their products in \(SS\) at the sum model
weight. The equality between \(\mathrm{holSub}\) and componentwise
\(SS\), followed by the source definition of
\(\mathrm{cuspW}\) and the already checked support, proves (IR9).
This is exactly
\srcid{ExistsUnramPtsDefinite.mul_mem_cuspOn_of_isSec}
followed by
\srcid{MulMemCuspFormsOfBoundedSectionComponent.mul_mem_cuspOn_of_bdd}.
\end{proof}

\paragraph{Frozen-source ledger and retained suppliers.}
All paths below use \srcid{alpoge-source-7cbf7a0c}.
In
\srcid{src/Development/AbelianVarietyTheory/Proofs/DivisorProofs/},
the actual reading proof is
\srcid{IsSecReadingDecreasingIntegerExponents.lean},
blob \srcid{b08455118b8cd54243ae8185790eaef6bf4e9382};
the final bounded component assembly is
\srcid{MulMemCuspFormsOfBoundedSectionComponent.lean},
blob \srcid{3a5a50359e5689723974601a420411cab829aa03}.
The earlier componentwise bounded-product argument is in
\srcid{src/Development/AbelianVarietyTheory/Definitions/Divisors/ExistsUnramPtsDefinite.lean},
blob \srcid{3cd59129266bb5e4fbc6a994794b40622e259d41}.
The actual nonvanishing supplier is
\srcid{src/Development/AbelianVarietyTheory/Proofs/IwasawaDecomposition/UnfoldedIsSecZeroKeySection.lean},
blob \srcid{54d610b59fda59d5c0fb5d17e846651747cf1cd4}.
Its level-\(3\) eta construction and the existence of the key unit
are separate actual inputs to the reading proposition.

The Fourier product proof is in
\srcid{src/ReciprocityTheory/ReductiveGroups/ArithmeticGeometry/FundamentalDomains/BoundedSectionAlgebra.lean},
blob \srcid{71f99c21a0d1bc7a47618506aea05641e5cad761}.
The concrete Fourier suppliers remain
\srcid{ModerateCoefficientPositivityHolds.lean}
(blob \srcid{712132fb6d8bf919371842e7a2a95d5f42d1fc57})
and
\srcid{CoefficientConvolutionHolds.lean}
(blob \srcid{dbd4d7e005c93fdc46adbd27159fb35b54cc2fbb})
in that same directory.
The latter itself uses
\srcid{CoefficientConvolutionOfSubCoefficients}
and \srcid{SublatticeCoefficients}.
This section proves the actual exponent clearing, local quotient,
and bounded component multiplication from those precise suppliers;
it does not substitute an unnamed Koecher principle for them.

\section{The source first-minimum estimate at a half maximum}
\label{sec:sr-source-half-maximum}
We keep the source frame datum, weighted norms, and rational
transvections throughout. The first-minimum estimate is made at
one actual frame of the same uniformization used by the theta reading.

\subsection{The periodic Schwarz estimate}
\begin{srlemma}\label{lem:sr-source-periodic-schwarz}
Let \(M>0\), \(0<t<T\), and let \(g:\mathbb C\to\mathbb C\) be
\(M\)-periodic, holomorphic on \(\operatorname{Im}z>-T\), and tend
to zero as \(\operatorname{Im}z\to+\infty\), uniformly in the real
coordinate in the sense of the source filter
\(\operatorname{comap}(\operatorname{Im})(+\infty)\).
If \(|g(z)|\le B\) on \(\operatorname{Im}z=-t\), then
\[
                    |g(0)|\le e^{-2\pi t/M}B .
 \tag{HM1}\label{eq:sr-source-periodic-schwarz}
\]
\end{srlemma}
\begin{proof}
Use \(q=e^{2\pi iz/M}\). Periodicity gives a single-valued
holomorphic function \(F(q)\) on the punctured disc of radius
\(e^{2\pi T/M}\).
The given limit makes \(F\) bounded near zero with \(F(0)=0\);
the removable singularity theorem, or its Cauchy coefficient proof,
makes this a holomorphic extension.
Set \(R=e^{2\pi t/M}>1\).
The divided difference \(F(q)/q\), with value \(F'(0)\) at zero,
is holomorphic on a neighborhood of the closed disc of radius \(R\).
On its boundary it has norm at most \(B/R\).
The maximum principle gives \(|F(1)|\le B/R\), which is (HM1).
This is precisely
\srcid{PeriodicSchwarzLemma.periodic_schwarz}.
\end{proof}

\subsection{Decay on the rational isotropic flow}
The source \(X=\mathrm{DataX}(E,r)\) includes the lattice \(L'\),
the scalar \(k_0>0\) with
\[
                     c_\tau=2|\tau(\xi_0)|=k_0,
\]
and an integer \(n_0>0\) sending \(L'\) into the lattice of
admissible rational isotropic vectors.
For \(v\in L'\) with \(v\ne0\) and
\(\langle v,v\rangle=0\), let \(x(\zeta)\) be the source flow
\(X.\mathrm{flow}(v,(\zeta)_\tau,x_0)\).

We use the exact source Fourier expansion on an admissible
isotropic direction and its coefficient-positivity statement.
The latter is the first Fourier input in
Section~\ref{sec:sr-source-integer-reading}.
For \(P\in SS_\lambda(X)\), write \(u=n_0v\).
The square in the rank-one flow matrix gives
\[
      x(\zeta)
        =X.\mathrm{flow}\bigl(u,(\zeta/n_0^2)_\tau,x_0\bigr).
 \tag{HM2}\label{eq:sr-source-flow-rescale}
\]
The source expansion is the convergent sum
\[
 P(x(\zeta))=\sum_\nu a_\nu
       \exp\!\left({4\pi i\zeta\over n_0^2}
                              \sum_\tau\nu_\tau\right)
                  \qquad(\operatorname{Im}\zeta\ge0),
 \tag{HM3}\label{eq:sr-source-isotropic-expansion}
\]
where \(a_\nu=X.\mathrm{fj}(\lambda,u,\nu,P,x_0)\).
Every nonzero \(a_\nu\) has all \(\nu_\tau\ge0\), and
\(a_0=0\) by the defining zero-coefficient condition of \(SS\).
Thus every nonzero term has \(\sum_\tau\nu_\tau>0\).

At \(\zeta=0\), the source's \(\mathrm{HasSum}\) in the finite
dimensional complex target gives \(\sum_\nu|a_\nu|<\infty\).
For \(\operatorname{Im}\zeta\ge0\), each summand in (HM3) has
norm at most \(|a_\nu|\), and its norm tends to zero as
\(\operatorname{Im}\zeta\to+\infty\).
Dominated convergence for this summable family therefore yields
\[
                         P(x(\zeta))\longrightarrow0 .
 \tag{HM4}\label{eq:sr-source-flow-decay}
\]
The argument applies to the source's function equal to
\(P(x(\zeta))\) on the flow domain and zero elsewhere, since its
values agree for all sufficiently large imaginary parts.
This is
\srcid{MaximumBoundField.flowTermNormDecay.tendsto_zero}.
It uses the actual Fourier expansion and positive-frequency
supplier; it does not assume rapid decay for an arbitrary
holomorphic function.

\subsection{The maximum estimate before host rescaling}
Set \(f=[E:\mathbb Q]/2\) and
\[
 \varepsilon_*=\min\!\left\{1,\,
 {\pi\over M(fw+2)k_0\log2}\right\}.
 \tag{HM5}\label{eq:sr-source-epstar}
\]
Assume \(M,w>0\), \(P\in SS_\lambda(X)\),
\(\sum_\tau\operatorname{spread}_\tau(\lambda)\le w\), and that
the source rational transvections are in the arithmetic group
with trivial character:
for every isotropic \(v\in L'\) and every integral
conjugation-fixed \(b\in E\), there is \(\gamma\in\Gamma_X\)
such that \(\chi_\lambda(\gamma)=1\) and
\[
                  \tau(\gamma)=\mathrm{flowMat}_{v,\tau}
                                                  (M\tau(b)).
 \tag{HM6}\label{eq:sr-source-flow-period}
\]
Let \(x_0\) be a positive half maximum:
\[
 0<\|P\|_{\lambda,x_0},\qquad
 \forall x\in\mathcal D_X:\quad
                 \|P\|_{\lambda,x}\le2\|P\|_{\lambda,x_0}.
 \tag{HM7}\label{eq:sr-source-half-max}
\]

\begin{srproposition}\label{prop:sr-source-maxbound}
For every \(0\ne v\in L'\), there is a good embedding \(\tau\)
with
\[
                 \varepsilon_*
                  \le\operatorname{maj}_{x_0,\tau}(\tau(v)).
 \tag{HM8}\label{eq:sr-source-maxbound}
\]
\end{srproposition}
\begin{proof}
Suppose every displayed majorant \(m_\tau\) is
less than \(\varepsilon_*\le1\).
The value \(\langle v,v\rangle\) is integral, and its every
complex conjugate has absolute value at most \(m_\tau<1\);
the other half of the embeddings are complex conjugates.
Its rational norm is therefore an integer of absolute value
less than one, so it vanishes.
The source majorant inequality thus makes \(v\) isotropic.

Each \(m_\tau>0\), since \(v\ne0\).
The explicit source flow-domain criterion is
\[
 x(\zeta)\in\mathcal D_X
      \quad\Longleftrightarrow\quad
                 1+k_0m_\tau\operatorname{Im}\zeta>0
                       \quad\text{for every }\tau .
 \tag{HM9}\label{eq:sr-source-flow-domain}
\]
Put \(T=(k_0\varepsilon_*)^{-1}\) and
\(t=(2k_0\varepsilon_*)^{-1}\), so \(0<t<T\).
The assumed strict inequalities imply that the flow lies
in the domain on \(\operatorname{Im}\zeta>-T\).
The flow matrix is polynomial in \(\zeta\), so
\(g(\zeta)=P(x(\zeta))\) is holomorphic there.
Taking \(b=1\) in (HM6), using the flow composition law and
the section's arithmetic transformation, makes \(g\)
\(M\)-periodic.
Equation (HM4) supplies its limit zero at imaginary infinity.

On the line \(\operatorname{Im}\zeta=-t\), all
\(1+k_0m_\tau\operatorname{Im}\zeta\) lie in \((1/2,1]\).
The source Gram-minor comparison for the isotropic flow gives
\[
 {n_\lambda(x_0)\over n_\lambda(x(\zeta))}
 \le\prod_\tau
   \max\{1+k_0m_\tau\operatorname{Im}\zeta,\,
             (1+k_0m_\tau\operatorname{Im}\zeta)^{-1}\}
                           ^{\operatorname{spread}_\tau(\lambda)/2}
 \le2^{w/2}.
\]
Together with (HM7), this implies
\[
     |g(\zeta)|\le
           {2\,\|P\|_{\lambda,x_0}\,2^{w/2}\over n_\lambda(x_0)}.
\]
Apply (HM1) and multiply by \(n_\lambda(x_0)>0\).
Since \(g(0)=P(x_0)\) and the half-maximum norm is positive,
\[
                  1\le e^{-2\pi t/M}\,2\,2^{w/2}.
 \tag{HM10}\label{eq:sr-source-maxbound-contradiction}
\]
But (HM5) gives
\(2\pi t/M=\pi/(Mk_0\varepsilon_*)\ge(fw+2)\log2\).
Since \(f\ge1\) and \(w>0\), the logarithm of the right side
of (HM10) is at most
\(-(fw+2)\log2+(1+w/2)\log2<0\).
This contradicts (HM10).
The proof is exactly \srcid{MaxBoundPf.maxBound_holds}.
The flow-domain and Gram-minor comparisons invoked above are the
source \srcid{FlowDomainCriterionHolds} and \srcid{M3bNormHolds}.
\end{proof}

\subsection{The host scalar and the root discriminant}
Write \(d_E=|\operatorname{disc}E|\) and
\(\mathrm{rd}_E=d_E^{1/[E:\mathbb Q]}\).
The principal host's two fractional ideals \(A,\bar A\) satisfy
\[
                 (\xi_0)A\bar A\,\mathfrak D_E=O_E,\qquad
                 |\tau(\xi_0)|=x>0\quad\text{for every }\tau .
 \tag{HM11}\label{eq:sr-source-host-ideal}
\]
Here \(\mathfrak D_E\) is the different.

\begin{srlemma}\label{lem:sr-source-small-inverse}
There is \(b\in E\), \(b\ne0\), such that
\(bA\subset O_E\) and
\[
                        |\tau(b)|^2\le x\,\mathrm{rd}_E^2
                            \quad\text{for every }\tau.
 \tag{HM12}\label{eq:sr-source-small-inverse}
\]
\end{srlemma}
\begin{proof}
First suppose \(N(\bar A)\le N(A)\).
Taking norms in (HM11), with \([E:\mathbb Q]=2f\), gives
\[
                  x^{2f}N(A)N(\bar A)d_E=1,
       \qquad N(A)^{-1}\le x^f\sqrt{d_E}.
\]
Set \(R=\sqrt{x\,\mathrm{rd}_E^2}\).
For the fractional ideal \(A^{-1}\), the complex-disc version
of Minkowski's theorem applies with this same radius at every
infinite place, because
\[
 2^fN(A)^{-1}\sqrt{d_E}
        \le2^fx^fd_E
        <\pi^fx^fd_E=\pi^fR^{2f}.
\]
It supplies \(0\ne b\in A^{-1}\) with every \(|\tau(b)|<R\).
This gives (HM12) and \(bA\subset O_E\).
If the norms have the opposite order, apply this argument with
\(A,\bar A\) reversed and conjugate \(b\).
The product identity is unchanged, conjugation exchanges the
ideals and permutes all embeddings, giving the same conclusion.
These are the source
\srcid{FirstMinAtHalfMax.exists_small_inv} and
\srcid{exists_small_inv_AA}.
\end{proof}

\subsection{The actual first minimum and the orthogonal half maximum}
Now fix the actual host packet \(D\), stage \(j\), supplied
uniformization laws \(\mathrm{Laws}(D,j,X_o,c,U,\mathrm{dict})\),
a good embedding \(\sigma\), and one component \(k\).
The source's principal-level construction gives a datum \(X\)
with this same underlying component, together with
\[
 \Lambda=A L',\quad k_0=2\,\xi_R(D),\quad
 M'=l.M\,N(\mathrm{ccK})\,D.\mathrm{host}.d,\quad
 M'k_0=2\,l.M\sqrt{D.\mathrm{host}.d},
 \tag{HM13}\label{eq:sr-source-transvection-data}
\]
where \(\Lambda=AL'\) means the actual \(O_E\)-span of
\(au\), \(a\in A,u\in L'\).
The transvections in (HM6) hold with \(M'\).
This precise supplier is
\srcid{ExistsSelfDualDataTransvectionPrincipalLevel.exists_dataX_transv}.
It uses the actual perfect trace form and the actual principal
congruence inclusion, not the abstract lattice alone.

Define
\[
 \epsilon_{\rm lev}(w)=
 \min\!\left\{1,\,
 {\pi\over2\,l.M\sqrt{D.\mathrm{host}.d}\,
                    (fw+2)\log2}\right\}.
 \tag{HM14}\label{eq:sr-source-epslev}
\]
Equation (HM13) identifies this with (HM5).
The source's \(\operatorname{firstMinimumSq}\) is the infimum,
over nonzero \(v\in\Lambda\), of
\(\sum_\tau 2\xi_R(D)\operatorname{maj}_{x_0,\tau}(\tau(v))\).

\begin{srproposition}\label{prop:sr-source-first-minimum-halfmax}
At a positive half maximum of a section with weight spread at
most \(w>0\) and trivial arithmetic character,
\[
 \operatorname{firstMinimumSq}(x_0)
             \ge {2\,\epsilon_{\rm lev}(w)\over\mathrm{rd}_E^2}.
 \tag{HM15}\label{eq:sr-source-first-minimum}
\]
\end{srproposition}
\begin{proof}
Choose \(b\) from (HM12) with \(x=\xi_R(D)\).
For \(v\in\Lambda\), the exact span in (HM13) and
\(bA\subset O_E\) give \(bv\in L'\).
For \(v\ne0\), (HM8) therefore supplies \(\tau_0\) with
\[
 \epsilon_{\rm lev}(w)
       \le\operatorname{maj}_{x_0,\tau_0}(\tau_0(bv))
       =|\tau_0(b)|^2\operatorname{maj}_{x_0,\tau_0}(\tau_0(v))
       \le\xi_R(D)\mathrm{rd}_E^2
                       \operatorname{maj}_{x_0,\tau_0}(\tau_0(v)).
\]
Every term in the first-minimum sum is nonnegative.
Keeping its \(\tau_0\)-term proves the claimed lower bound for
every nonzero \(v\).
The lattice spans \(E^r\), and the actual packet has \(r\ge1\),
so the set defining the infimum is nonempty.
Taking the infimum proves (HM15).
This is \srcid{FirstMinAtHalfMax.firstMin_at_halfMax}.
\end{proof}

\begin{srproposition}\label{prop:sr-source-orthogonal-halfmax}
Suppose a family \(P\) belongs to the actual
\(\mathrm{holSub}(U_\sigma,\lambda)\), has positive finite
global supremum \(S\), has trivial arithmetic characters, and
every component has spread at most the same \(w>0\).
There are a component \(k\) and an orthogonal frame \(x_f\) in
its domain such that
\[
              S\le2|P_k(x_f)|,\qquad
 \operatorname{firstMinimumSq}(x_f)
             \ge {2\,\epsilon_{\rm lev}(w)\over\mathrm{rd}_E^2}.
 \tag{HM16}\label{eq:sr-source-orthogonal-halfmax}
\]
\end{srproposition}
\begin{proof}
Since \(0<S<\infty\), the supremum definition gives a frame
\(x\) in one component with \(S/2<\|P_k\|_{\lambda,x}\).
All component norms are at most \(S\), so \(x\) is a positive
half maximum in that component.
The component belongs to \(SS\), because the source
\(\mathrm{holSub}\) is the span of such families and \(SS\)
is a submodule.
Apply (HM15).
Use the actual upper triangular orthogonalizing factor \(b(x)\)
from \srcid{FrameDegreeData.bOf_mem}.
Right multiplication by it preserves the frame domain, the
weighted norm, and \(\operatorname{firstMinimumSq}\).
At the resulting orthogonal frame all successive Gram-minor
norms are one; its weighted norm is therefore \(|P_k(x_f)|\).
This gives both inequalities in (HM16).
It is \srcid{FirstMinAtHalfMax.half_sup_at_ortho}.
No maximum-attainment or compactness assertion is used.
\end{proof}

\paragraph{Frozen-source ledger and scope.}
All sources are read at \srcid{alpoge-source-7cbf7a0c}.
The periodic Schwarz proof is
\srcid{src/ReciprocityTheory/ReductiveGroups/ArithmeticGeometry/FundamentalDomains/PeriodicSchwarzLemma.lean},
blob \srcid{0f49e1e5cf5bd6b20b73373433d9600b5f02d11f}.
In that same directory,
\srcid{MaximumBoundField.lean} is
\srcid{c7b65dca30b0f3b9a858814b7e5bf9e5f42bb774},
\srcid{MaxBoundPf.lean} is
\srcid{dcd344008b85a949193a31cd613956653dddeb08}, and
\srcid{CuspidalExponentialDecay.lean} is
\srcid{a05dc166cf19232be5e2c25ba52b9a7713d53d3e}.
The host-scalar, first-minimum and half-sup proof is
\srcid{src/Development/AbelianVarietyTheory/Proofs/DivisorProofs/FirstMinAtHalfMax.lean},
blob \srcid{63e29ac9732ebfa986c8251e1b2f6917e907e707}.
The exact principal-level supplier (HM13) is
\srcid{ExistsSelfDualDataTransvectionPrincipalLevel.lean} in
that directory, blob \srcid{8a14c6e08bbc4bcc1218f16298c4ed44e1b165d4}.

The source Fourier expansion \srcid{M3bExpHolds}, coefficient
positivity, flow-domain formula \srcid{FlowDomainCriterionHolds},
and Gram-minor comparison \srcid{M3bNormHolds} are retained as
their exact analytic suppliers in the foregoing proof.
The actual perfect trace-form and principal-congruence
construction in (HM13), and the same uniformization's Borel
orthogonalization identities in (HM16), are likewise the
specified source inputs. The pointwise theorem does not assert
a bound uniform over independently chosen metrics or frames.

\section{The source Gaussian supremum and derivative estimates}
\label{sec:sr-source-gaussian-sup}
We use the source's Euclidean normalization
\(H(z,w)=\sum_i z_i\overline{w_i}\), Lebesgue measure on
\(\mathbb C^N\), and Gaussian \(e^{-q\pi\|z\|^2}\).
The square of the \(L^2\)-norm is denoted \(L_2\), as in the source.

\subsection{The actual lattice and its covolume}
Let \(Y=\mathrm{GT}(\iota)\), \(N=|\iota|\).
Its period lattice \(L\subset\mathbb C^N\) is generated by a
real basis \(b_1,\ldots,b_{2N}\).
The alternating form \(\operatorname{Im}H\) is integral on \(L\),
and every \(\mathbb Z\)-linear functional on \(L\) is pairing
with a point of \(L\). The supplied factor \(k(a)\) has
\(|k(a)|=1\) for \(a\in L\).
An actual level-\(q\) section is an entire \(F\) satisfying
\[
 F(z+a)=k(a)\exp\{q\pi H(z,a)+q\pi H(a,a)/2\}F(z).
 \tag{GS1}\label{eq:sr-source-GT-section}
\]
Its norm density and the source normalized integral are
\[
 h_F(z)=|F(z)|^2e^{-q\pi\|z\|^2},\qquad
 L_2(F)={1\over\operatorname{vol}D}\int_Dh_F,\qquad
 \lambda_1=\inf_{0\ne a\in L}\|a\|,
 \tag{GS2}\label{eq:sr-source-GT-norm}
\]
where \(D\) is the half-open basis parallelepiped.
Equation (GS1) gives \(h_F(z+a)=h_F(z)\) directly.

The source first proves \(\operatorname{vol}D\le1\).
Let \(B\) be the real basis matrix in orthonormal coordinates
and \(J\) the orthogonal complex structure representing
\(\operatorname{Im}H\), with \(|\det J|=1\).
The alternating Gram matrix is \(G=B^tJ^tB\).
For each lattice coordinate functional, perfection supplies a
lattice vector representing it.
Writing these vectors in the given integral basis produces
an integer matrix \(C\) with \(CG=I\).
Thus \(\det C\ne0\), \(|\det C|\ge1\), and
\[
                  1=|\det C|\,|\det B|^2,\qquad
                           \operatorname{vol}D=|\det B|\le1 .
 \tag{GS3}\label{eq:sr-source-GT-volume}
\]
This follows the source
\srcid{LemmaSA.lemSG_inv}, \srcid{lemS_detM_le},
and \srcid{lemS_vol_le}.
In particular the integral over \(D\) is at most \(L_2(F)\).
For \(N>0\), discreteness of the basis lattice and the existence
of a nonzero basis vector give \(0<\lambda_1<\infty\).

\subsection{The local Gaussian mean}
For an entire \(G:\mathbb C^N\to\mathbb C\), \(N>0\),
\[
 |G(0)|^2\le {2\over\pi^N}
       \int_{\|z\|\le\sqrt{2N}}|G(z)|^2e^{-\|z\|^2}\,dz .
 \tag{GS4}\label{eq:sr-source-gaussian-ball-mean}
\]
Here is the source's mean argument.
For fixed \(z\), the one-variable Cauchy mean and
Cauchy--Schwarz give
\(|G(0)|^2\le\int_{|u|=1}|G(uz)|^2\,du\), with probability
measure on the circle.
Integrating on a radial ball and using the measure-preserving
rotations \(z\mapsto uz\) proves
\[
 |G(0)|^2\int_{\|z\|\le\sqrt{2N}}e^{-\|z\|^2}
       \le\int_{\|z\|\le\sqrt{2N}}|G(z)|^2e^{-\|z\|^2}.
\]
The total Gaussian mass is \(\pi^N\), and its second moment is
\(N\pi^N\). On \(\|z\|>\sqrt{2N}\),
\(1\le\|z\|^2/(2N)\), so that region has mass at most
\(\pi^N/2\).
This proves (GS4).
The source verifies these two moments through the radial
integrals
\(\int_0^\infty r^ke^{-r^2}dr=\Gamma((k+1)/2)/2\).
For \(N=0\) the ball has its single point with mass \(1\);
(GS4) remains true.
These are the steps in \srcid{RSubMean.sq_norm_zero_le_gauss_ball}.

For \(q\ge1\), \(v_0\in\mathbb C^N\), put
\(s=\sqrt{q\pi}\), \(w=z/s\), and
\[
 G(z)=F(v_0+w)\exp\{-q\pi H(w,v_0)\}.
\]
Then
\[
 |G(z)|^2e^{-\|z\|^2}
              =e^{q\pi\|v_0\|^2}h_F(v_0+w).
\]
Changing real \(2N\)-dimensional volume by \(s^{-2N}\) in
(GS4) gives the exact local estimate
\[
 h_F(v_0)\le2q^N
       \int_{\|u-v_0\|\le\sqrt{2N}/\sqrt{q\pi}}h_F(u)\,du .
 \tag{GS5}\label{eq:sr-source-local-Gaussian}
\]
This is \srcid{LemmaSB.lemS_local}, with its actual square-root
and \(\pi\) normalization.

\subsection{Packing and the global estimate}
A finite \(\delta\)-separated subset of a ball of radius \(\rho\)
in \(\mathbb C^N\) has cardinality at most
\[
                         (1+4\rho/\delta)^{2N}.
 \tag{GS6}\label{eq:sr-source-complex-packing}
\]
To match the source's constants, when \(\rho<\delta/2\)
there is at most one point.
Otherwise place disjoint open balls of radius \(\delta/2\)
around the points and compare their real \(2N\)-dimensional
volumes with the containing ball of radius
\(\rho+\delta/2\). This gives
\((1+2\rho/\delta)^{2N}\); the source's
\(\mathrm{PackingCount}\) bound \( (3\max(\delta,\rho)/\delta)^{2N}\)
also implies (GS6) in the stated range.
This is the conversion in \srcid{LemmaS.lemS_pack}.

Let \(\rho=\sqrt{2N}/\sqrt{q\pi}\) and
\(K=(1+4\rho/\lambda_1)^{2N}\).
For any fixed \(u\in D\), the lattice translates \(u+a\) lying
in the closed ball of radius \(\rho\) around \(v_0\) form a
\(\lambda_1\)-separated set.
Apply (GS6) to each finite subset; the nonnegative sum of their
indicators is therefore at most \(K\).
Unfolding Lebesgue measure over \(D+L\), Tonelli and the
periodicity of \(h_F\) yield
\[
 \int_{\|u-v_0\|\le\rho}h_F(u)\,du
   \le K\int_Dh_F(u)\,du\le K L_2(F).
 \tag{GS7}\label{eq:sr-source-ball-cover}
\]
The density is continuous and the closure of \(D\) is compact,
so the latter integral is finite.
This justifies the conversion from the nonnegative integrals,
including the finite-measure requirements.
For \(N=0\) the same estimate is the one-point calculation.

\begin{srproposition}\label{prop:sr-source-gaussian-sup}
For every actual \(F\in Y.\mathrm{Sec}(q)\), \(q\ge1\),
\[
 |F(v)|^2e^{-q\pi\|v\|^2}
 \le 2q^N
       \left(1+{4\sqrt{2N}\over\sqrt{q\pi}\lambda_1}\right)^{2N}
                     L_2(F)\qquad(v\in\mathbb C^N).
 \tag{GS8}\label{eq:sr-source-gaussian-sup}
\]
For \(N=0\), the parenthesized factor is interpreted by the
separate zero-dimensional calculation.
\end{srproposition}
\begin{proof}
Combine (GS5) and (GS7).
This is exactly \srcid{LemmaS.lemS_main}, then
\srcid{gSup} and \srcid{GlobalTheory.sup_le_l2_principal_torus}.
\end{proof}

\subsection{The source repeated Cauchy estimate}
Let \(F\) be entire on \(\mathbb C^N\), and suppose
\[
                   |F(z)|^2e^{-\pi\|z\|^2}\le S
                                      \quad\text{for every }z.
\]
Necessarily \(S\ge0\).
For a list of \(d\) coordinate directions,
\[
 |D_{i_1}\cdots D_{i_d}F(0)|
                   \le\sqrt S\,(\pi e d)^{d/2}.
 \tag{GS9}\label{eq:sr-source-Gaussian-Cauchy}
\]
For \(d=0\) the last factor is \(1\).
The source proves this by repeated one-variable Cauchy estimates.
If \(|F|\le A\) whenever \(\sum_i|z_i-x_i|\le d\rho\), induction
on the list gives the derivative bound \(A/\rho^d\).
In the induction step apply the one-variable Cauchy derivative
bound on the circle of radius \(\rho\) in the next coordinate;
each remaining ball stays in the original \(d\rho\)-ball.

For \(d>0\), take \(\rho=(\pi d)^{-1/2}\).
On that \(d\rho\)-ball about zero,
\(\|z\|^2\le(\sum_i|z_i|)^2\le d/\pi\).
The Gaussian hypothesis gives \(|F(z)|\le\sqrt S\,e^{d/2}\).
Consequently the repeated bound is
\(\sqrt S\,e^{d/2}(\pi d)^{d/2}\), which is (GS9).
This is the actual \srcid{RCauchy.iter_step_bound},
\srcid{cauchy_gauss_le} and \srcid{GCauchy.gCauchy} argument.
For a homogeneous polynomial \(P\) of degree \(d\), finite
linearity gives
\[
 |P(\partial)F(0)|
     \le \left(\sum_\alpha|P_\alpha|\right)
                           \sqrt S\,(\pi e d)^{d/2}.
 \tag{GS10}\label{eq:sr-source-polynomial-Cauchy}
\]
The source polynomial derivative is the sum over monomials of
the iterated derivative along their actual repeated coordinate
list; no factorial or divided-power convention is inserted.

\subsection{The exact self-dual short-vector supplier}
Let \(E\) be CM, \(f=[E:\mathbb Q]/2\), and let
\(\Lambda\subset E^n\) be a finitely generated \(O_E\)-module
self-dual for \(h(u,v)=\sum_k u_k\overline{v_k}\).
The source's rational positive form is
\[
 B(u,v)=\operatorname{Tr}_{E/\mathbb Q}h(u,v),\qquad
 B(u,u)=\sum_\sigma\sum_k|\sigma(u_k)|^2 .
 \tag{GS11a}\label{eq:sr-source-trace-hermitian}
\]
It is symmetric, positive definite and nondegenerate.
Since \(\Lambda\) is full over \(E\), it is full over \(\mathbb Q\):
clear an integer denominator in each scalar \(a\in E\), so
\(a\Lambda\subset\mathbb Q\Lambda\).
A \(\mathbb Z\)-basis of \(\Lambda\) is therefore a
\(\mathbb Q\)-basis \(b_1,\ldots,b_M\) of \(E^n\), \(M=2fn\).

For a nonzero vector \(u\) in the \(B\)-dual of this lattice,
\(B(u,y)\in\mathbb Z\) for all \(y\in\Lambda\).
Writing \(\mathfrak D\) for the different, testing at
\(\overline a\,y\), \(a\in O_E\), gives
\(h(u,y)\in\mathfrak D^{-1}\).
Hence \(\delta u\in\Lambda\) for every
\(\delta\in\mathfrak D\), by Hermitian self-duality.
Put \(\alpha=h(u,u)\ne0\).
For \(\delta\in\mathfrak D\) and \(a\in O_E\),
\(a\delta u\in\Lambda\); the trace integrality just obtained,
together with invariance of trace under conjugation, gives
\(\operatorname{Tr}(\alpha\delta a)\in\mathbb Z\).
Thus
\[
                  \alpha\mathfrak D\subset\mathfrak D^{-1},
                    \qquad \alpha\in\mathfrak D^{-2}.
\]
The norm of a nonzero element of a fractional ideal is at least
the ideal norm, since its principal ideal divided by that
fractional ideal is integral. Consequently
\(|N_{E/\mathbb Q}\alpha|\ge|\operatorname{disc}E|^{-2}\).
All conjugates of \(\alpha\) are nonnegative, with positive sum.
The arithmetic--geometric mean inequality gives
\[
                   B(u,u)\ge {2f\over\mathrm{rd}_E^2}.
 \tag{GS11b}\label{eq:sr-source-trace-dual-min}
\]
This is the original
\srcid{EmbeddingHermSelf.trace_dual_min}, using the fractional
ideal norm and trace estimate already proved by the trace-dual
norm lemma.

Apply the source transference bound to this rational Euclidean
space with \(\lambda=\sqrt{2f}/\mathrm{rd}_E\).
For clarity, its rational-to-real passage diagonalizes \(B\)
and embeds the rational space into \(\mathbb R^M\) by multiplying
the basis coordinates by the square roots of its nonnegative
diagonal coefficients.
The images of the \(B\)-dual basis form a real basis:
pairing with the images of the original basis gives the identity
matrix. Their integer lattice has minimum at least \(\lambda\)
by (GS11b). The original Euclidean transference theorem supplies
an independent family in the dual lattice with lengths at most
\(M^2/\lambda\). Its integral pairings with that real basis are
integer coordinates in the original rational basis, so this
family pulls back to \(v_1,\ldots,v_M\in\Lambda\), independent
and spanning over \(\mathbb Q\), with
\(B(v_j,v_j)\le(M^2/\lambda)^2\).
This is \srcid{TraceHermNondegenerate.transference_rat}, using
the earlier source \srcid{T0.transference_dual} inequality.

For every embedding \(\tau\), its conjugate is distinct and has
the same squared-coordinate sum, so
\(2\sum_k|\tau(v_{j,k})|^2\le B(v_j,v_j)\).
Substitution of \(M=2fn\) and \(\lambda=\sqrt{2f}/\mathrm{rd}_E\)
proves
\[
 \sum_k|\tau(v_{j,k})|^2
       \le\{2 f^{3/2}n^2\mathrm{rd}_E\}^2 .
 \tag{GS11}\label{eq:sr-source-tensor-short-vectors}
\]
This proves the actual supplier
\srcid{TraceHermNondegenerate.exists_short_vectors}, including
its stronger rational independence and spanning assertion.
For \(n=0\) the empty family is exactly the asserted family.

\paragraph{The tensor-period minimum used with this estimate.}
Let \(T\) be the source \(O_E\)-linear positive torus, and
\(\Lambda\subset E^n\) the actual positive self-dual Hermitian
coefficient lattice.
The preceding source short-vector construction gives an
\(E\)-spanning family \(v_j\in\Lambda\) with (GS11).
For \(\gamma\) in the actual tensor lattice, contraction by
\(v_j\) lies in the original period lattice.
If \(\gamma\ne0\), some such contraction is nonzero:
otherwise linearity and the spanning property make all
contractions zero, hence every coordinate of \(\gamma\) zero.
Cauchy--Schwarz gives
\[
 \|\operatorname{contr}_{v_j}\gamma\|
     \le2 f^{3/2}n^2\mathrm{rd}_E\,\|\operatorname{flat}\gamma\|.
\]
Thus, if every nonzero original period has length at least
\(\lambda>0\), every nonzero tensor period has length at least
\[
                  \lambda'={\lambda\over
                                     2f^{3/2}n^2\mathrm{rd}_E}.
 \tag{GS12}\label{eq:sr-source-tensor-period-minimum}
\]
For \(n=0\) the tensor lattice is zero and the source takes
\(\lambda'=1\).
These are the exact contraction steps in
\srcid{LeMulNormFlatOfSelfDual.le_mul_norm_flat_of_selfDual}
and \srcid{NormalisedThetaValueBound.exists_min_period}.
The short-vector construction (GS11) has just been proved from
the same source transference bound and trace-dual norm estimate.

\paragraph{Frozen-source ledger.}
All sources use \srcid{alpoge-source-7cbf7a0c}.
In
\srcid{src/ReciprocityTheory/ReductiveGroups/ArithmeticGeometry/SquareIntegrability/},
the relevant files and blobs are
\begin{itemize}
\item \srcid{GCauchyStatement.lean},
\srcid{e53576b249c74c71fcb3af8c66b5d3111df34e04};
\item \srcid{LemmaSA.lean},
\srcid{dbdc66e846263d475c90d6c0c9ef0e1c580b845a};
\item \srcid{RSubMean.lean},
\srcid{d6dda1e32c137194deb3cf921537bf7c21e6efe4};
\item \srcid{LemmaSB.lean},
\srcid{948ae7e2d3bc2acbcf2b3a97ad9b4ceae6fe1609};
\item \srcid{LemmaS.lean},
\srcid{1950e4726b3b5ce278491239f53fcab26d875218};
\item \srcid{RCauchy.lean},
\srcid{07fa9a67f8ac8ff589a18c3691ec38b76f963408};
\item \srcid{GCauchy.lean},
\srcid{22dc8c28fccd0c9fad68fd47686cd6e0a01cef93}.
\end{itemize}
The realization of the packing supplier by disjoint ball volumes
is \srcid{src/Development/AbelianVarietyTheory/Proofs/CuspFormEvaluations/SupLeL2PrincipalTorus.lean},
blob \srcid{89e09ea66f30995d745920d4a3f50db3921618ba}.
The tensor contraction proof in the same directory is
\srcid{LeMulNormFlatOfSelfDual.lean},
blob \srcid{0800466bdb169dc1ddd430812943ca9c4f8f093d};
the actual normalized tensor-period bound is
\srcid{NormalisedThetaValueBound.lean},
blob \srcid{689a12045c213d4d11e73f6660d5ee8a8739f684}.
The exact short-vector proof in the same directory is
\srcid{TraceHermNondegenerate.lean},
blob \srcid{c12f943d2da5d109d91c74539874b853f035eb6b},
with \srcid{EmbeddingHermSelf.lean},
blob \srcid{d1edfe1b816a8069a1c99f902b09cbf9f6549fd1}.
Its fractional-ideal trace and norm supplier is
\srcid{src/Development/AbelianVarietyTheory/Proofs/CovolumeBounds/FinrankMulDiscrRpowLeSumSqOfMemTraceDual.lean},
blob \srcid{83942a82c35de9c2751c6b1419702f69ca8ab550}.
Its Euclidean transference supplier is the already identified
\srcid{src/ReciprocityTheory/ReductiveGroups/ArithmeticGeometry/HeightTheory/Transference.lean},
blob \srcid{3b1ced809264736e3140deb8ea437d8aa061580c}.
The Gaussian statements above are proved for a supplied actual
\(\mathrm{GT}\); they do not construct the model's analytic
fibre or identify an algebraic theta reading with its sections.

\section{Projected vacua and the raw-theta index factor}
\label{sec:sr-source-raw-theta-norm}
This section follows the source normalized isogeny sum, followed
by \srcid{JetTranslateProjectedVacuumLe} and
\srcid{NormRawThetaSeriesIndexBoundLeSqrt}.
It keeps the same Hermitian ambient space, actual Weil action,
frame and algebraic theta reading in every comparison.

\subsection{The normalized isogeny sum on two actual tori}
Use two source Gaussian tori \(Y,Y'\) on the same
\(\mathbb C^N\), with lattices \(L,L'\), unit factors \(k,k'\),
and the actual translation
\[
 (U_aF)(z)=
       e^{\pi H(z,a)-\pi H(a,a)/2}F(z-a).
 \tag{RV1}\label{eq:sr-source-fock-translation}
\]
The factors obey their integral semicharacter law and agree on
\(L\cap L'\).
Let \(R\subset L'\) be a finite set of unique representatives
of \(L'/(L\cap L')\), and set
\[
                     G_R=\sum_{a\in R}k'(a)U_aF
                                      \qquad(F\in Y.\mathrm{Sec}(1)).
 \tag{RV2}\label{eq:sr-source-isogeny-sum}
\]

\begin{srlemma}\label{lem:sr-source-isogeny-sum-norm}
The actual sum \(G_R\) belongs to \(Y'.\mathrm{Sec}(1)\), and
\[
                           L_{2,Y'}(G_R)=|R|\,L_{2,Y}(F).
 \tag{RV3}\label{eq:sr-source-isogeny-sum-norm}
\]
\end{srlemma}
\begin{proof}
The Fock composition law and the semicharacter law say that
\(k'(a)U_a\) is an additive representation for \(a\in L'\).
For \(\gamma\in L'\), write
\(\gamma+a=a'+b\), uniquely with \(a'\in R\) and
\(b\in L\cap L'\).
The operator \(k'(b)U_b=k(b)U_b\) fixes \(F\).
Thus translation by \(\gamma\) permutes the summands in (RV2),
which is the section condition for \(Y'\).
Each summand is entire.

The finite representative set supplies a full common sublattice
\(L_0\subset L\cap L'\): for each real basis vector \(b_i'\)
of \(L'\), two integer multiples have the same representative
modulo \(L\), so a nonzero integer multiple \(n_i b_i'\)
belongs to \(L\); these multiples are still a real basis.
Integrate on its basis parallelepiped.
For \(a,a'\in R\), put
\[
 h_{a,a'}(z)=U_aF(z)\overline{U_{a'}F(z)}e^{-\pi H(z,z)} .
\]
Substitution of (RV1) and the section condition gives exactly
\[
 h_{a,a'}(z+x)
   =e^{2\pi i\,\operatorname{Im}H(x,a-a')}h_{a,a'}(z)
                                             \qquad(x\in L).
 \tag{RV4}\label{eq:sr-source-pairing-density}
\]
For \(x\in L_0\) this multiplier is one, by integrality in \(L'\).
For \(a\ne a'\), their difference is not in \(L\).
Perfection of \(L\) therefore gives some \(x\in L\) for which
the multiplier in (RV4) is not one.
Indeed if every such imaginary pairing were integral,
the resulting integral functional would be represented by
a point of \(L\); nondegeneracy over the real span would then
put \(a-a'\) in \(L\).
The integral of the \(L_0\)-periodic continuous density is
unchanged by any translation of its fundamental domain.
Equation (RV4) consequently forces its off-diagonal integral
to be zero.

On the diagonal,
\(h_{a,a}(z)=|F(z-a)|^2e^{-\pi\|z-a\|^2}\).
Its average over a fundamental domain of \(L_0\) equals
\(L_{2,Y}(F)\), independently of \(a\), by tiling with
fundamental domains of \(L\).
Similarly the average of the density of \(G_R\), periodic
under \(L'\), equals \(L_{2,Y'}(G_R)\).
Expanding the finite square, all off-diagonal averages vanish
and each of the \(|R|\) diagonal terms has the indicated value.
The factors \(k'(a)\) have modulus one.
All integrals are legitimate because the densities are continuous
on bounded fundamental parallelepipeds.
This proves (RV3), by exactly the common-lattice and pairing-density
argument in
\srcid{ContinuousThetaSectionPointwiseNormSq.isogeny_sum_mem_and_norm}.
\end{proof}

\subsection{The normalized tensor jet}
Let \(T\) be the source \(O_E\)-linear positive torus, with
level-one semicharacter determined by the fixed integral \(w_0\)
satisfying \(w_0+\overline{w_0}=1\).
Let \(s\in T.\mathrm{Sec}(1,1)\), \(s\ne0\).
For a positive self-dual Hermitian coefficient lattice
\(\Lambda\subset E^n\), use its actual tensor lattice and
representatives \(R\) modulo the original product period lattice.
Both tensor tori carry the same formula
\[
 \epsilon_{w_0}(\gamma)
 =\exp\{\pi i\,\operatorname{Im}
                         H_C(\gamma,w_0\gamma)\}.
\]
They therefore have the required equality of factors on their
intersection.

The source obtains a real basis of this tensor lattice by
integer commensurability with the product lattice.
Integrality and perfection are transported by its contraction
description: contraction against every coefficient vector in
\(\Lambda\) lies in the original period lattice.
The Hermitian self-duality of \(\Lambda\) supplies a finite dual
family, which reconstructs a tensor from these contractions.
Thus the tensor object used here is the actual
\(\mathrm{GT}\), not merely a collection of unrelated norm bounds.

For \(F=s^{\otimes n}\), the product fundamental domain and
Fubini give \(L_2(F)=L_2(s)^n\).
Lemma~\ref{lem:sr-source-isogeny-sum-norm} gives
\[
                         L_2(G_R)=|R|\,L_2(s)^n .
 \tag{RV5}\label{eq:sr-source-tensor-sum-norm}
\]
Also \(L_2(s)>0\).
If it were zero, its continuous nonnegative norm density would
vanish almost everywhere on one fundamental domain, hence
everywhere there and, by the actual section transformation,
everywhere; this would imply \(s=0\).

Put \(f=[E:\mathbb Q]/2\), \(g=\dim_\mathbb C T=fr\).
Suppose the original period lattice has minimum at least
\(\lambda>0\).
The exact short-vector and contraction argument in
Section~\ref{sec:sr-source-gaussian-sup} gives the tensor
minimum
\(\lambda'=\lambda/(2f^{3/2}n^2\mathrm{rd}_E)\) for \(n>0\);
the zero-row case is handled separately as there.
For a homogeneous polynomial \(P\) of degree \(d\), let
\(\|P\|_1=\sum_\alpha|P_\alpha|\), and define
\[
 {\cal B}(P,d,\lambda)=
 \|P\|_1(\pi e d)^{d/2}\sqrt2
 \left(1+
 {4\sqrt{2gn}\,(2f^{3/2}n^2\mathrm{rd}_E)
                       \over\sqrt\pi\,\lambda}\right)^{gn}.
 \tag{RV6}\label{eq:sr-source-normalized-jet-bound}
\]
The degree-zero factor is \(1\).

Apply the actual Gaussian supremum bound to \(G_R\) on this
tensor torus, then the repeated Cauchy estimate to any translate
\(U_yG_R\). The norm density of \(U_yG_R\) at \(z\) is the
norm density of \(G_R\) at \(z-y\), so its same global bound
applies. For coefficients \(|c_j|\le1\), a finite set \(J\),
and \(|\beta|^2=|R|\), the result is
\[
 { \left|\beta^{-1}\sum_{j\in J}c_j
                         P(\partial)U_{y_j}G_R(0)\right|
                       \over L_2(s)^{n/2}}
                  \le |J|\,{\cal B}(P,d,\lambda).
 \tag{RV7}\label{eq:sr-source-normalized-isogeny-jet}
\]
Here \(\beta\ne0\), since \(R\) contains a representative of zero.
The cancellation is literal:
the square root of (RV5) is
\(|\beta|L_2(s)^{n/2}\).
This is
\srcid{NormalisedThetaValueBound.normalised_theta_value_bound}.
No beta factor is chosen independently of the actual index.

\subsection{The actual projected vacuum at an orthogonal frame}
Keep the source Hermitian datum \(X\), one fixed symplectic
rational basis \(\mathrm{bas}\), and the corresponding genuine
Weil and Heisenberg actions.
Let \(L\subset E^r\) be self-dual for the source alternating
trace form \(EF_X\), \(n\) even, and let \(v\) be the normalized
vacuum for \(\mathrm{tens}(O_E^n,L)\).
For the positive coefficient lattice \(\Lambda\), put
\[
 J_\Lambda=[\Lambda:\Lambda\cap O_E^n],\qquad
                         A_\Lambda=\mathrm{tens}(\Lambda,L).
\]
The source tensor-index formula gives a finite representative set
\(R\subset A_\Lambda\) modulo \(\mathrm{tens}(O_E^n,L)\), with
\[
                  |R|=J_\Lambda^r=:J>0,\qquad
 \mathrm{projA}_{A_\Lambda}v
       ={1\over J}\sum_{a\in R}\epsilon_{w_0}(a)\,\omega(a)v .
 \tag{RV8}\label{eq:sr-source-projector-representatives}
\]
To identify this with the definition as the mean over distinct
orbit functions, use the source vacuum orthogonality:
\(\epsilon(a)\omega(a)v=v\) holds exactly on its original
self-dual lattice. The Heisenberg composition law then gives
equality of two orbit functions exactly when their parameters
differ by that lattice.
Representatives therefore map bijectively to the distinct
orbit functions, including their count \(J\).
This is \srcid{ProjAEqSumReps.projA_eq_sum_reps} and
\srcid{vac_proj_mean}.
The tensor-index formula itself is the source
\srcid{TensorLatticeSelfDualSemicharacterIndex}:
a pseudobasis of either factor makes the tensor relative index
the product of the factor indices; fractional ideal multiplication
does not change each relative index.
It gives
\(\operatorname{relIndex}(\Lambda'\otimes L,\Lambda\otimes L)
 =\operatorname{relIndex}(\Lambda',\Lambda)^r\).
Self-duality gives the required finite nonzero index and vacuum
overlap. We retain the source tensor-vacuum construction used
in this step.

Fix an orthogonal frame \(x\in\mathcal D_X\), and a common
\(t>0\) with \(X.c_\tau=t\) for every \(\tau\).
The source constructs the positive torus \(T\) and
\(\psi:\mathbb C^g\overset{\sim}\longrightarrow V_T\)
from this same \(X,L,x\).
In the actual relabelled coordinates,
\[
 H_T(\psi z,\psi z')=h_x(z,z'),\qquad
 \psi(e_j)=\sqrt t\,e_{\epsilon(j)},\qquad
 \psi(p_v(x,u))\in L_T\ \Longleftrightarrow\ u\in L.
 \tag{RV9}\label{eq:sr-source-self-torus}
\]
For precision, this construction first uses the full lattice
\(L\), an integer commensurability with \(X.\Lambda\), and a
real lattice basis of its scaled period image.
The integral dual of \(L\) represents every lattice functional:
extend it rationally and invert the nondegenerate \(EF_X\)
pairing. Self-duality puts the representing vector back in \(L\).
The torus factor is the sign of
\(\operatorname{Im}H_T(\gamma,w_0\gamma)\);
the integer cross-term formula and \(w_0+\overline{w_0}=1\)
give its semicharacter law.
Finally \(\psi\) is the explicit coordinate permutation followed
by multiplication by \(\sqrt t\).
Uniqueness of the Hermitian form with the given imaginary
pairing on a full period lattice identifies it with \(h_x\).
These are the steps in
\srcid{ExistsTorusOfSelfD.exists_torus_of_selfD}.
The equality of the period coordinates with the source \(h_x\)
uses the specified \srcid{HermitianFormAtFrameLaws}.

Let the actual normalized theta function satisfy
\[
 \theta\in\mathrm{IsTheta1}(X,x,L),\qquad
 \mathrm{MeanOne}(\theta),\qquad
       \prod_b\theta(z_b)=\kappa\,\mathrm{thA}(v,x,z).
 \tag{RV10}\label{eq:sr-source-theta-vacuum-tie}
\]
Under (RV9), \(s(v_1)=\theta(\psi^{-1}v_1)\) is a level-one
section with \(L_2(s)=1\).
Indeed substitution in its Fock transformation gives (GS1),
and the normalized integral is unchanged by the fixed
real linear coordinate change: its Haar scalar cancels between
the integral and volume. Any two fundamental domains give the
same integral of the periodic density.
Thus \(\mathrm{MeanOne}\) is exactly this \(L_2\)-normalization;
it is not an unrelated positive scalar.

Write \(\Phi(y)_b=\psi(p_v(x,y_b))\).
The source proves that \(-\Phi(R)\) is the representative set
for the tensor torus above,
\(\epsilon_{w_0}(-\Phi(a))=\epsilon_{w_0}(a)\),
and Fock translation becomes \(U_{-\Phi(y)}\).
Consequently the theta function for
\(\omega(y)\mathrm{projA}_{A_\Lambda}v\) is
\[
                  \kappa^{-1}J^{-1}\,G_R\circ\psi^{\,n},
\]
with the additional common translation \(U_{-\Phi(y)}\).
If \(\kappa=0\), the target bound is immediate.
Otherwise the displayed \(\kappa^{-1}\) cancels its exterior
\(\kappa\).

Each actual theta-polynomial letter becomes
\(\sqrt t\) times its Euclidean coordinate direction.
The relabelling of letters is injective, so it preserves the
coefficient \(L^1\)-sum.
Homogeneity gives the derivative factor \(t^{d/2}\).
Apply (RV7) with one translation and \(\beta=\sqrt J\).
Since (RV8) divides by \(J\), the exact result is
\[
 |\kappa\,\mathrm{thP}(\omega(y)\mathrm{projA}_{A_\Lambda}v,P,x)|
           \sqrt J
 \le \|P\|_1(\pi e t d)^{d/2}\sqrt2
 \left(1+
 {4\sqrt{2gn}\,(2f^{3/2}n^2\mathrm{rd}_E)
                       \over\sqrt\pi\,\lambda}\right)^{gn}.
 \tag{RV11}\label{eq:sr-source-projected-vacuum-jet}
\]
Here the minimum assumption is exactly
\(\lambda^2\le h_x(p_v(x,u),p_v(x,u)).\operatorname{re}\)
for every \(0\ne u\in L\).
The source \srcid{hx_ortho} supplies this from the same frame's
\(\lambda^2\le\mathrm{firstMinimumSq}(x)\).
Equation (RV11) is
\srcid{JetTranslateProjectedVacuumLe.jet_projA_le}.

\subsection{The exact raw-theta value}
Now use the original \(E,H,D_w,l,X_o,c\), good embedding
\(\sigma:c.F\to\mathbb C\), actual \(\mathrm{IsUnf}\) with
\(r=p+q\ge1\), and actual \(\mathrm{IsNab}\) connection.
The row dictionary in the preceding construction has
\(X.p_0=D_w.q\) and the other block size \(q_0=D_w.p\), exactly
as in \srcid{ExistsAnalyticDataUnitFrameFramedPoint.exists_setup}.
The polynomial directions, signatures and \(\psi\) use that
specified dictionary, including this ordering.

Let \(n\) be even, \(a\) dominant, and let
\(\Lambda=\mathrm{La}\) be the actual positive coefficient
lattice with \(\mathrm{Unimod}\),
\(\mathrm{TrivAt}(l.M)\) and \(\mathrm{LevOK}(N,l.M)\).
Let \(P\in\mathrm{harmL}(c,\Lambda,a)\) have total degree
\(d=\mathrm{DF}(D_w,a)\).
Let \(f\in\mathrm{sec}(X_o,c,a,n)\), and suppose its values on
every source test ring and every frame satisfy the actual
\(\mathrm{thLaValAt}\) formula for these same
\(\Lambda,N,P,\mathrm{nab}\) and coefficient function \(w\).
Assume \(|\sigma(w(q))|\le1\) for every label \(q\).

The source first extends the equality from test rings to the
complex frame: actual \(\mathrm{secBaseChange}\) supplies a
test frame mapping to it; uniqueness and naturality of the
actual short theta value give the same value after pullback.
This is \srcid{thLaValAt_evR}.
It uses all test rings in the premise, not equality at one
preferred complex point.

At an orthogonal frame \(x\), the supplied analytic fibre,
marked level data, theta generator and normalized vacuum give
the exact complex value formula. If
\[
 m=\mathrm{mIdx}(\Lambda,N,r),\quad
 z=\mathrm{zCard}(\Lambda,N,r),\quad
 Q=|\mathrm{Fin}(n)\to\mathrm{LevQ}(D_w,l)|,\quad
 J=J_\Lambda^r,
\]
its scalar is
\[
 {z\over m}\ {1\over Q}\,
 \kappa\,\mathrm{thP}
       \left(\sum_q\sigma(w(q))\,
                  \omega(Y_q)\mathrm{projA}_{A_\Lambda}v,\ P,\ x\right),
              \qquad z=mJ .
 \tag{RV12}\label{eq:sr-source-raw-reading-scalar}
\]
The exact relation \(z=mJ\) follows by multiplying the finite
lattice indices through the intersection with \(O_E^n\);
it is \srcid{zCard_eq}, with \srcid{index_Ln}.
No normalizing \(\mathrm{betaH}^{-1}\) is inserted in this raw
formula.

Put \(B_0\) equal to the right side of (RV11).
Linearity of the actual theta derivative and the triangle
inequality give a bound \(B_0/\sqrt J\) for the normalized
\(Q\)-mean in (RV12).
Combining this with \(z=mJ\) gives
\[
 \left|\operatorname{ev}_\sigma(f;x)\right|
 \le \sqrt{J_\Lambda^r}\,
       \left(\sum_\alpha|\sigma(P_\alpha)|\right)
       (\pi e t d)^{d/2}\sqrt2
 \left(1+
 {4\sqrt{2frn}\,(2f^{3/2}n^2\mathrm{rd}_E)
                       \over\sqrt\pi\,\lambda}\right)^{frn}.
 \tag{RV13}\label{eq:sr-source-raw-value-bound}
\]
Algebraically \(|m^{-1}z|\le J\) even in the source's
zero-denominator convention; the actual finite-index case has
\(m>0\).
Likewise an empty-label mean is zero by that convention.
The actual label group is finite and nonempty.
This proves exactly
\srcid{NormRawThetaSeriesIndexBoundLeSqrt.norm_evR_le}
under the stated actual analytic-fibre/reading suppliers.

\paragraph{Frozen-source ledger and retained geometric comparisons.}
All sources use \srcid{alpoge-source-7cbf7a0c}.
In
\srcid{src/Development/AbelianVarietyTheory/Proofs/CuspFormEvaluations/},
the files are
\begin{itemize}
\item \srcid{ContinuousThetaSectionPointwiseNormSq.lean},
\srcid{9f7ab3413f449c9b93dac642ca2ab0845057ee85};
\item \srcid{NormalisedThetaValueBound.lean},
\srcid{689a12045c213d4d11e73f6660d5ee8a8739f684};
\item \srcid{ProjAEqSumReps.lean},
\srcid{818dccb76df27daa5ee0a1c63a15a27700ebfaf4};
\item \srcid{ExistsTorusOfSelfD.lean},
\srcid{859af98c532c58c34785cd7f6e382a0f49944c71};
\item \srcid{L2OneEqIntegralFundamentalDomainDivVolume.lean},
\srcid{34d7c34bb82c87d1f3b0a8a9ea0771ffa8c0d23e};
\item \srcid{JetTranslateProjectedVacuumLe.lean},
\srcid{4a358a4e34143451dab2d499357ed33d1e51249d};
\item \srcid{ExistsAnalyticDataUnitFrameFramedPoint.lean},
\srcid{903ac3995f8acf7789604f184c54eb7fc7887e45}.
\end{itemize}
The tensor lattice and relative-index proof is
\srcid{src/Development/AbelianVarietyTheory/Proofs/SiegelThetaTransformation/TensorLatticeSelfDualSemicharacterIndex.lean},
blob \srcid{4894f884d42f8c17e88ea2604bcf2fc40817b20d}.
The final raw reading proof is
\srcid{src/Development/AbelianVarietyTheory/Proofs/DivisorProofs/NormRawThetaSeriesIndexBoundLeSqrt.lean},
blob \srcid{d1817b6beba51a259498dc4ae71cb669cf67c874}.

The remaining construction-specific suppliers at (RV10)--(RV12)
are exactly the actual
\srcid{AnalyticFibreUnitFrameUniformisation},
\srcid{FramedPointGeneratorIsThetaLevelOne},
\srcid{ProdLevelOneThetaEqVacuumTheta},
\srcid{DualHalfLevelMarkingExistsAndMatch},
and
\srcid{ShortThetaValueComplexPointEqClassicalTheta},
together with actual \srcid{IsNab}, the symplectic row dictionary,
and the source tensor-vacuum/Heisenberg laws.
The preceding arguments prove the Gaussian and index bounds with
these compatible data; they do not construct separate witnesses
and identify them merely because they have the same names.
